\documentclass[10pt,slovak,a4paper]{article}

\usepackage[slovak]{babel}
\usepackage[IL2]{fontenc}
\usepackage[utf8]{inputenc}
\usepackage{graphicx}
\usepackage{url} 
\usepackage{hyperref}
\usepackage{tabularx}
\usepackage{changepage}
\usepackage{float}
\usepackage{xcolor, colortbl}
\usepackage{booktabs}

\definecolor{Gray}{rgb}{0.88,0.87,0.88}
\newcolumntype{a}{>{\columncolor{Gray}}c}
\newcolumntype{e}{>{\columncolor{Gray}}l}
\newcolumntype{f}{>{\columncolor{Gray}}r}
\newcolumntype{d}{>{\columncolor{Gray}}p{0.4\textwidth}}
\newcolumntype{z}{>{\columncolor{Gray}}X}

\newcommand{\thinline}{\arrayrulecolor{gray}\specialrule{0.2pt}{0pt}{0pt}\arrayrulecolor{black}}


\renewcommand{\arraystretch}{1.4}
\setlength{\extrarowheight}{2pt}
\begin{document}


\section*{Záväzná osnova pre aplikovaný výskum a vývoj}
Maximálny rozsah záväznej osnovy je 15 strán

\begin{table}[H]
\centering
\begin{adjustwidth}{-2cm}{-2cm}
\begin{tabularx}{\linewidth}{|l|X|}
\hline
Názov projektu & Personalizovaný asistent sústredenia študentov \\ \hline
Zodpovedný riešiteľ & Martina Bublincová  \\ \hline
Žiadateľ & Fakulta informatiky a informačných technológií STU v Bratislave  \\ \hline
Štatutárny zástupca žiadateľa &   \\ \hline

\end{tabularx}
\end{adjustwidth}
\end{table}

\section{Excelentnosť}

\begin{itemize}
\item Charakterizujte projektový zámer, navrhovaný spôsob riešenia a aplikačnú úroveň
očakávaných výsledkov
\item Definujte aktuálnosť cieľov projektu pre danú aplikačnú oblasť vedy a techniky
a vysvetlite v čom spočíva originalita myšlienky projektu a originalita jeho riešenia
\item Charakterizujte stav poznania, vedecké základy a predchádzajúce výsledky,
ktoré sú základom pre realizáciu projektu
\item Definujte ciele projektu a reálnosť ich dosiahnutia
\item Opíšte navrhovanú metodológiu riešenia projektu, opodstatnenosť jej výberu vo vzťahu
k splneniu deklarovaných cieľov
\item Charakterizujte maximálne 5 najdôležitejších výstupov aplikačného výskumu
zodpovedného riešiteľa za posledných 5 rokov; uveďte ich dôležitosť na národnej
a medzinárodnej úrovni
\item Charakterizujte maximálne 3 najdôležitejšie projekty realizované zodpovedným
riešiteľom/-kou za posledných 5 rokov v štruktúre informácií: názov projektu, grantová
schéma, roky riešenia, rozpočet projektu, pozícia zodpovedného riešiteľa v projekte
(hlavný riešiteľ / riešiteľ), objasnite význam výstupov projektu v národnom
aj medzinárodnom kontexte
\item Charakterizujte úroveň kvality zodpovedného riešiteľa v aplikačnej oblasti výskumu,
jeho expertízu a víziu pre budovanie výskumného tímu, ďalší rozvoj výskumu a kariérny
rast
\item Opíšte kompetentnosť zúčastnených riešiteľských a realizačných organizácií
na riešenie predkladaného projektu
\item Opíšte expertízu jednotlivých riešiteľov na riešenie predkladaného projektu
(netýka sa zodpovedného riešiteľa)
\item Definujte mieru zapojenia mladých pracovníkov výskumu a vývoja do 35 rokov vrátane
študentov doktorandského štúdia do riešenia projektu
\end{itemize}
 
Zámerom projektu je pomôcť s pozornosťou pri učení študentov.\\
Navrhovaný spôsob riešenia je pomocou mobilnej aplikácie, určenej na zber, \\
spracovanie a vizualizáciu údajov o rušení pozornosti počas štúdia a následnú pomoc s limitáciu rušiacich aplikácií.\\
Očakávame, že aplikácia pomôže zvýšiť pozornosť pri štúdiu.  \\

V dnešnej dobe majú študenti čím ďaľej tým väčšie problémy s pozornosťou, ako aj zmienené v našich zdrojoch.\\
Tým, čím sa líšime od ostatných riešení je úplný zámer na štúdium a študentov.  \\

Čerpáme zdroje z viacerých publikácií, ktoré poukazujú na vplyv mobilných notifikácií/dištraktorov na pozornosť študentov.  \\

Náš cieľ projektu je mať pozitívny vplyv na pozornosť študentov počas štúdia.\\
Existujú, iné aplikácie ktoré úspešne riešia viacej generálnu problematiku blokovaní aplikácií,
a teda naše riešenie zamerané na študentov by malo mať ešte lepší efekt vďaka špecializácii. \\

Riešitelia sú študenti Slovenskej Technickej Univerzity na Fakulte Informatiky a Informačných Technológií.\\

\section{Dopad}
\begin{itemize}
\item Opíšte prínosy projektu na rozvoj poznania, aplikovaného výskumu a inovácií
(nové produkty alebo postupy, zdokonalené produkty alebo postupy, nové úžitkové
materiály, nové technológie, nové služby, spoločenské inovácie, a pod.)
\item Opíšte spôsob a mieru využiteľnosti deklarovaných výsledkov riešenia projektu
pre žiadateľa/užívateľa v SR a v zahraničí
\item Definujte mieru ekonomického prínosu (zvýšenie podielu pridanej hodnoty, úspora
zdrojov, dopad na zamestnanosť, využitie domácich surovín, skvalitnenie ľudských
zdrojov, zlepšenie kvality života a/alebo zdravia, ochrana životného prostredia)
\item Definujte opatrenia na maximalizáciu výsledkov a dopadu ako aj plánovanú
komunikáciu projektových výstupov

\end{itemize}

Daný projekt by mal pomôcť spresniť negatívny dopad rušiacich aplikácií na štúdium a mitigáciu tohto dopadu.\\
Užívateľ si stiahne aplikáciu a pomocou danej aplikácie môže reálne ovplyvniť kvalitu štúdia. \\
Zlepší to znalosti a kvalitu budúcej pracovnej sily.\\
Budeme aktívne testovať vplyv aplikácie počas všetkých etáp a na základe spätnej väzby upravovať aplikáciu aby mala najväčší dopad.  \\
Anonymizované výsledky budeme zdieľať s akademickým svetom.  

\clearpage
\section{Implementácia}
\begin{itemize}
\item Projektový plán a výstupy - definujte realizáciu projektu z hľadiska časového
harmonogramu s ohľadom na stanovené ciele projektu, popíšte jednotlivé kroky
pracovného plánu s jednotlivými míľnikmi, rozdeľte ich do pracovných balíkov
a definujte a opíšte plánované výstupy z riešenia projektu
\item Opíšte projektový manažment a procesy v rámci riadenia projektu
\item Opíšte riziká implementácie projektu vo vzťahu k očakávaným výstupom a spôsob ich
zmiernenia
\item Vysvetlite adekvátnosť navrhnutého rozpočtu projektu v kontexte finančnej náročnosti
dosiahnutia cieľov
\item Opíšte existujúcu technickú a personálnu infraštruktúru pracoviska (pracovísk)
podieľajúcich sa na implementácii projektu

\end{itemize}
Realizácia projektu začne s vytvorením prototypu aplikácie. \\
Tento prototyp bude následne testovaný dva týždne aspoň 30 študentmi.  \\
Zozbierané dáta a spätná väzba budú použité na zlepšenie prototypu počas mesiaca v etape 3.  \\
V etape 4 prichádza druhá vlna testovania taktiež o dĺžke 2 týždne.  \\
V poslednej etape príde na záverečné vylepšenia, vytvorenie finálneho produktu a zozbieranie anonymizovaných dát pre akademický svet.  \\

Časové plánovanie projektu bude riadené v Jire a všetky dôležité súbory budú na tímovom repozitári na Githube.  \\
Každý kritický bod projektu bude predchádzaný tímovým stretnutím.  \\
Hlavné riziká implementácie sú ignorovanie všetkých návrhov našej aplikácie počas štúdia a strach zo zneužitia osobných dát.  \\
Mitigácia dopadu ignorovania návrhov počas testovania bude pomocou finančnej odmeny a výsledky testovania sa použijú aj na zamedzenie tohto rizika.  \\
Strach zo zneužitia osobných dát bude zmiernený tým, že spravíme repozitár verejným a všetky dáta posielané na server budú anonymizované.  \\

Daný rozpočet zahŕňa ceny testovania a vydania aplikácie do obchodov Google Play a App Store a server na Hetzneri.\\
Inak budú financie potrebné len na samotný vývoj aplikácie a prípadne dokončíme bez platu finálnu verziu.  \\

Pracovať budeme na vlastných zariadeniach napr. laptopy, telefóny, počítače atď. \\
Server na zbieranie dát budeme mať na Hetzneri.  \\

\end{document}
